% ======================================================================
\section{Results}\label{sec:results}

All numbers are raw counts from the scorer or values derived from those counts with Eq.~\ref{eq:cost}
and~\ref{eq:costw}; derived values are marked. The counts of the final system are reproduced exactly by an
independent re-export of every decision (Section~\ref{sec:repro}). Appendix~\ref{app:sources} gives the result file
behind each number.

\subsection{Main result}\label{sec:headline}
\begin{table}[h]
\centering\small
\caption{Main result. Recall = hits / needed; precision = hits / (hits + wrong); ``show nothing'' costs
$4\cdot\text{needed}/\text{clips}$ (all three derived).}
\label{tab:headline}
\resizebox{\linewidth}{!}{%
\begin{tabular}{lrrrrrrrr}
\toprule
 & \multicolumn{4}{c}{development set (71 clips, 58 needed)} & \multicolumn{4}{c}{test set (88 clips, 65 needed)} \\
\cmidrule(lr){2-5}\cmidrule(lr){6-9}
system & hits & wrong (v/c/p) & cost & recall / prec. & hits & wrong (v/c/p) & cost & recall / prec. \\
\midrule
show nothing & 0 & 0 & 3.268 & -- & 0 & 0 & 2.955 & -- \\
September baseline & 18 & 51 (5/37/9) & 3.690 & .31 / .26 & 21 & 40 (2/31/7) & 2.909 & .32 / .34 \\
\textbf{final system} & \textbf{29} & \textbf{15 (6/7/2)} & \textbf{2.056} & \textbf{.50 / .66} & \textbf{24} & \textbf{24 (4/15/5)} & \textbf{2.409} & \textbf{.37 / .50} \\
\bottomrule
\end{tabular}}

\smallskip
\footnotesize v / c / p = wrong pictures of an on-screen sound / of another sound / with nothing there.
\end{table}

\begin{itemize}
  \item On the development set the final system finds half of the needed sounds, and two of three of its pictures are
        right. On the test set it finds 37\,\% and half of its pictures are right.
  \item The September baseline was \emph{worse than showing nothing} on the development set (3.690 against 3.268, and
        still 3.521 with a longer repeat merge) and only slightly better on the test set (2.909 against 2.955). The final
        system is clearly cheaper than showing nothing on both (2.056 and 2.409).
  \item Against the baseline, the final system has 11 more hits and 36 fewer wrong pictures on the development set, and
        3 more hits and 16 fewer wrong pictures on the test set (cost $-1.634$ and $-0.500$ per clip). About a third of
        the development gain carries over to the test set; since many variants were compared on the development set,
        part of its gain is expected to be selection luck, and the test numbers are the better estimate for new videos.
  \item On the test set most wrong pictures belong to another sound (15 of 24); pictures of on-screen sounds are few
        (4).
\end{itemize}

\paragraph{Significance.} The paired bootstrap of the final system against the September baseline on the test set has
not been run: it needs per-clip rows that are stored on the cluster, and the baseline's rows were scored with a
shorter repeat merge (2.0\,s instead of 2.5\,s), which raises its wrong-picture count. An intermediate version of the
final system, without its last display, grouping and witness rules, was cheaper than the baseline on the test set by
$-0.341\ [-0.636, -0.045]$ per clip (one-sided $p = 0.017$, about 0.03 two-sided); the final system has a lower cost
than that version on both sets. Appendix~\ref{app:history} gives every intermediate comparison.

\subsection{Sensitivity}\label{sec:sensitivity}
\paragraph{On-screen weight.} With Eq.~\ref{eq:costw} the final system costs 1.887 / 1.972 / 2.056 on the development
set and 2.318 / 2.364 / 2.409 on the test set at $w = 0 / 1 / 2$ (derived). For every accepted change from the
baseline to the final system the cost falls at every $w$, so no decision depends on how on-screen pictures are
weighted.

\paragraph{Cost weights.} The weights 4 and 2 were fixed before the final development phase and not varied for the
final system. On the September test table (Section~\ref{sec:sept-table}) the gated system was the cheapest of all
systems for any wrong-picture weight $\beta$ between 0.39 and 2.56, which contains $\beta = 2$.

\subsection{Where the misses come from}\label{sec:errors}
For every missed needed sound, the logged decision trail records the step at which it was lost:

\begin{table}[h]
\centering\small
\caption{Fate of the missed needed sounds of the final system.}
\label{tab:misses}
\begin{tabular}{L{4.0cm}rrL{8.4cm}}
\toprule
 & dev. & test & where \\
\midrule
misses & 29 & 41 & \\
\quad lost at a named step & 20 & 32 & development: listener band rescue 6, on-screen gate 5, DASM vote 2, family merge 2, mirror veto 2, scene check 1, masked-weak veto 1, one rescue per family 1. Test: band rescue 6, continuation veto 5, family merge 4, gate 3, mirror veto 3, DASM vote 2, BEATs extraction 2, seven other steps 1 each \\
\quad never heard by any detector & 6 & 6 & no span of the family near the onset \\
\quad timing / matching & 3 & 3 & a same-family span exists but starts outside the window \\
\bottomrule
\end{tabular}
\end{table}

Most misses are not ``never heard'': the sound was a candidate and a later rule removed it, most often because the
listeners did not confirm a weak detection. The on-screen gate silences few needed sounds (5 and 3), but it is also why
the on-screen wrong pictures remain (6 and 4).

\subsection{How far could this pipeline go?}\label{sec:ceiling}
A ceiling analysis replaces one stage at a time by the gold answer on the \emph{same} detections. It was run on the
saved pictures of an earlier version (25 of 55 needed sounds, 27 wrong, cost 2.451 on the development set; the
final system adds display and witness rules on top), so its absolute numbers belong to that version; its proportions
are the best available map.

\begin{table}[h]
\centering\small
\caption{Single-stage oracles (development set, earlier version, 55 needed).}
\label{tab:ceiling}
\begin{tabular}{lrrrr}
\toprule
oracle (one stage made perfect) & hits & wrong (v/c/p) & cost & $\Delta$ \\
\midrule
none (version as run) & 25 & 27 (9/16/2) & 2.451 & -- \\
perfect on-screen gate & 30 & 18 (0/16/2) & 1.915 & $-0.536$ \\
perfect listener, filters lifted & 34 & 24 (7/15/2) & 1.859 & $-0.592$ \\
perfect vetoes / filters & 30 & 27 (9/16/2) & 2.169 & $-0.282$ \\
perfect timing & 26 & 20 (11/7/2) & 2.197 & $-0.254$ \\
perfect family & 25 & 26 (14/10/2) & 2.423 & $-0.028$ \\
\midrule
all stages perfect (ceiling) & 43 & 3 (0/1/2) & 0.761 & $-1.690$ \\
\bottomrule
\end{tabular}
\end{table}

Applied in sequence, the gate gives the largest single step ($-0.676$). Even with every stage perfect, 12 of 55 needed
sounds stay missed: 5 are heard by no model at all (a hammer, an explosion, a clang, two golf whacks) and 7 are too
faint or too badly timed. This is the \emph{detector floor}: with today's detectors, no decision layer can find more than
about four in five needed sounds.

\subsection{What the on-screen gate is worth}\label{sec:sept-table}
The value of the gate itself was measured against a \emph{blind} system (same detector and picture model, no gate) on
the September test table (60 test clips, the baseline's detector, the same gate model and rule):

\begin{table}[h]
\centering\small
\caption{Gated minus blind, September test table; two-sided $p$ with Holm correction.}
\label{tab:sept}
\begin{tabular}{lrrl}
\toprule
measure & gated $-$ blind & 95\,\% interval & Holm $p$ \\
\midrule
F1 (pre-registered primary) & $+0.059$ & $[-0.030,\ +0.144]$ & 0.18 (unadjusted, null) \\
precision & $+0.137$ & & 0.010 \\
wrong pictures per clip & $-0.517$ & $[-0.800,\ -0.283]$ & $<0.001$ \\
viewer cost per clip & $-0.833$ & & 0.012 \\
clean-clip accuracy & $+0.219$ & & 0.006 \\
\midrule
unseen clips: gated vs show nothing & $-3.077$ & & $<0.001$ \\
all clips: gated vs show nothing & $-0.233$ & & 0.585 \\
\bottomrule
\end{tabular}
\end{table}

The gate removes about half a wrong picture per clip at the price of a small loss in recall. Its value lies almost
entirely in what it does \emph{not} show, which is why F1 -- which weighs a miss and a wrong picture equally -- does not
move. This comparison was not repeated with the final detector.

\subsection{Other measurements}\label{sec:other-results}
\paragraph{Gate accuracy.} On the labelled development sounds the gate has balanced accuracy 0.617: it silences 16 of
43 on-screen sounds and keeps 31 of 36 needed ones. A smaller model of the same family (Qwen2.5-VL-7B) reached 0.61.
Shifting the sampled frames by 0.5\,s changes 6 of 79 verdicts (7.6\,\%).

\paragraph{Listener agreement as evidence.} On the screening set, a detection that neither listener names is a real
sound 22\,\% of the time, one named by one listener 37\,\%, and one named by both 62\,\%. This is the basis of the witness
rule. A listener asked about the whole clip proposes extra sounds that are real 80\,\% of the time (74 of 92), against
38--39\,\% for the raw detectors; they still become wrong pictures in the full pipeline, because the gate passes
off-screen ones whose name is wrong, or because the labels list only the salient sounds.

\paragraph{Onset timing.} On 255 DCASE 2025 Task 3 events, the onset method places 80\,\% of the detected events within
0.5\,s (mean absolute error 0.35\,s). Only 26\,\% of those short-clip events reach the display bar at all, which again
points to detection, not timing, as the limit.

\paragraph{Pictures.} In a blind glance test by the author (54 sounds, 1.5\,s per picture, no name shown), Qwen-Image
pictures were recognised 26 times against 14 for the earlier FLUX.1 pictures with the same text
($+0.222\ [+0.074, +0.370]$). One of the 26 sent a false message (a telephone bell drawn as a desk bell).
